\section{DM--Electron Scattering (QEDark / QCDark / QCDark2)}
\label{sec:dme}

\subsection{Physical process and Lagrangian}

A sub-GeV DM particle $\chidm$ scatters off a bound valence electron in the
target crystal, promoting it to the conduction band and depositing a
continuous recoil-energy spectrum. The benchmark interaction (dark-photon
mediator; \citealt{Essig2012,Essig2016}) is
\begin{equation}
\mathcal{L} \supset -g_\chi\,\bar{\chidm}\gamma^\mu\chidm\,A'_\mu
- e\,Q_e\,\bar{e}\gamma^\mu e\,A_\mu
- \frac{\kappa}{2}\,F^{A'}_{\mu\nu}F^{\mu\nu}\,.
\end{equation}
After kinetic-mixing diagonalization this generates an effective
DM--electron vertex with coupling $g_\chi g_e \kappa / m_\mathrm{med}^2$,
where $m_\mathrm{med}$ is the mediator mass. Two limits define the two
``mediator cases'' used throughout CCDarkSens:
\begin{itemize}[leftmargin=*]
  \item \textbf{Heavy mediator} ($m_\mathrm{med} \gg$ typical momentum
        transfer $q \sim \alpha m_e$): the propagator is momentum-independent,
        $\FDM(q) = 1$.
  \item \textbf{Light (ultralight/massless) mediator} ($m_\mathrm{med} \ll q$):
        the propagator $\propto 1/q^2$, giving $\FDM(q) = (\alpha m_e / q)^2$.
\end{itemize}

\subsection{Reference cross section and mediator form factor}

The free-electron DM cross section, evaluated at the reference momentum
transfer $q_0 = \alpha m_e$ (the typical momentum transfer in atomic-scale
scattering), is
\begin{equation}
\bar{\sigmae} = \frac{\mue^2\, g_\chi^2\, g_e^2}{\pi\, m_{A'}^4}\bigg|_{q=q_0}\,,
\end{equation}
where $\mue = m_\chidm m_e/(m_\chidm + m_e)$ is the DM--electron reduced
mass. This single number, $\sigmae$, is the scan parameter reported on the
$y$-axis of the final limit. The momentum-transfer dependence factors out
into the dimensionless mediator form factor:
\begin{equation}
\FDM(q) = \begin{cases}
1 & \text{heavy } (n_\mathrm{FDM}=0) \\
(\alpha m_e / q)^2 & \text{light } (n_\mathrm{FDM}=2)
\end{cases}
\end{equation}
CCDarkSens's Python rate modules
(\texttt{ccdarkphys.common.mediator\_map.MEDIATOR\_TO\_FDM\_INDEX}) encode
this exactly as the QEDark convention $\FDM(q) \propto q^{-n_\mathrm{FDM}}$,
$n_\mathrm{FDM} \in \{0,2\}$.

\subsection{The crystal form factor}

The electronic structure of the target enters through the \textbf{crystal
form factor} $f_\mathrm{crystal}(\mathbf{q},\omega)$, which encodes the
probability of a momentum-$q$, energy-$\omega$ transition from an occupied
valence Bloch state to an empty conduction state:
\begin{equation}
|f_\mathrm{crystal}(\mathbf{q},\omega)|^2 \propto \frac{1}{V_\mathrm{cell}}
\sum_{i\to f} |\langle f | e^{i\mathbf{q}\cdot\mathbf{r}} | i\rangle|^2\,
\delta(E_f - E_i - \omega)\,.
\end{equation}
In the random-phase approximation (RPA) this is exactly proportional to the
imaginary part of the inverse dielectric function at finite momentum:
\begin{equation}
|f_\mathrm{crystal}(\mathbf{q},\omega)|^2 \;\propto\;
\Im\!\left[-\frac{1}{\varepsilon(\mathbf{q},\omega)}\right].
\end{equation}
CCDarkSens has two independent implementations of this quantity, both
consumed downstream by the identical C++ \texttt{RateTable}/
\texttt{ccdarksens\_scan\_dmelectron\_pattern} pipeline:

\begin{center}
\begin{tabular}{@{}p{2.1cm}p{6.5cm}p{5cm}@{}}
\toprule
Engine & Source of $|f_\mathrm{crystal}|^2$ or $\varepsilon(q,\omega)$ & Files \\
\midrule
QEDark & Precomputed $|f_\mathrm{crystal}|^2$ lookup table
  (\texttt{Si\_f2.txt}, from \citealt{Essig2016} DFT on bulk Si) &
  \texttt{python/ccdarkphys/qedark/entry.py} \\
QCDark2 & Full $\varepsilon(q,\omega)$ HDF5 from DFT+RPA (in-house
  dielectric-function code), on a $(q,E)$ grid with local-field effects &
  \texttt{python/ccdarkphys/qcdark2/entry.py},
  \texttt{data/qcdark2\_epsilon/<material>/*.h5} \\
\bottomrule
\end{tabular}
\end{center}

QCDark2 is the more general engine: it accepts an arbitrary DFT-derived
$\varepsilon(q,\omega)$ HDF5 (attributes \texttt{M\_cell}, \texttt{V\_cell},
\texttt{dE}), so the same rate-generation code path handles standard Si and
any other material for which a dielectric function file exists
(Section~\ref{sec:dme-hypmat}).

\subsection{DM velocity distribution and kinematics}

The local DM population follows the truncated Standard Halo Model (SHM):
\begin{equation}
f(\mathbf{v}) = \frac{1}{N_\mathrm{esc}}
\exp\!\left(-\frac{|\mathbf{v}+\mathbf{v}_E|^2}{v_0^2}\right)
\Theta(\vesc-|\mathbf{v}+\mathbf{v}_E|)\,.
\end{equation}
CCDarkSens's production configs
(\texttt{configs/examples/qedark\_generate\_*.json},
\texttt{qcdark2\_generate\_*.json}) use
\begin{equation}
v_0 = 238\ \mathrm{km/s}, \qquad v_E = 263\ \mathrm{km/s}, \qquad
\vesc = 544\ \mathrm{km/s}\,.
\end{equation}
For a deposited $(q,\omega)$, the minimum DM speed able to supply that
energy and momentum transfer is
\begin{equation}
\vmin(q,\omega) = \frac{\omega}{q} + \frac{q}{2m_\chidm}\,,
\end{equation}
and the halo integral reduces to the mean inverse speed above threshold:
\begin{equation}
g(\vmin) = \int_{v>\vmin} \frac{f(\mathbf{v})}{v}\,d^3v\,.
\end{equation}

\subsection{The $dR/dE$ master formula}

Combining the previous three subsections, the differential DM--electron
scattering rate per unit target mass per unit deposited energy is
\begin{equation}
\boxed{
\frac{dR}{dE}\bigg|_\omega = \frac{\rhochi}{\rhoT\, m_\chidm}\,\bar{\sigmae}
\int \frac{d^3q}{(2\pi)^3}\, \FDM^2(q)\,\frac{1}{q}\,\frac{2\pi^2}{\omega}\,
\Im\!\left[-\frac{1}{\varepsilon(q,\omega)}\right]\, g(\vmin)
}
\label{eq:dRdE_dme}
\end{equation}
with $\rhochi = 0.3$~GeV/cm$^3$ the local DM density and $\rhoT$ the target
mass density. This is evaluated numerically by both QEDark and QCDark2,
producing a CSV of $dR/dE$ [events/(kg$\cdot$yr$\cdot$eV)] versus $E$ for
each $(m_\chidm, \sigmae, \text{mediator})$ grid point, in exactly the
two-column format read by \texttt{RateTable::LoadCSV} and rebinned onto a
uniform ROOT histogram by \texttt{RateTable::MakeTH1D}
(\texttt{src/io/RateTable.cc}), which \textbf{linearly interpolates the
tabulated rate at each output bin's center}. For DM-electron and Migdal
spectra — smooth and continuous over many bins — bin-center sampling is an
excellent approximation. It is \emph{not} safe for the dark photon
channel's monochromatic boxcar signal; see
Section~\ref{sec:binning} for why, and for the numerical fix (\texttt{nbins})
required there.

\textbf{Threshold behavior:} $\varepsilon(q,\omega)$ (and hence $dR/dE$) is
zero for $\omega$ below the material's band gap $\Egap$ — this is where the
target's band structure sets the absolute low-mass reach of the search.

\subsection{Standard silicon case}

\begin{center}
\begin{tabular}{@{}p{4.3cm}p{4cm}p{5.5cm}@{}}
\toprule
Parameter & Value & Note \\
\midrule
Band gap $\Egap$ & 1.2~eV (QEDark/scissor) or 1.11~eV (darkelf optical, 130~K) &
  Both values appear in this codebase for different purposes (see Section~\ref{sec:klein}) \\
Mean $e$-$h$ pair energy $\ehpair$ & 3.8~eV & Measured, 100~K Si, PRD 102, 063026 \\
Target density $\rhoT$ & 2.329~g/cm$^3$ & Bulk Si \\
Valence electron count & $Z_\mathrm{val}=4$ (Si, $3s^2 3p^2$) & Sets number of scattering targets per unit cell \\
Crystal form factor source & \texttt{Si\_f2.txt} (QEDark) or \texttt{Si\_comp.h5} (QCDark2) & Both reproduce the same physics; QCDark2 includes local-field effects \\
$\sigmae$ reference $q_0$ & $\alpha m_e \approx 3.7$~keV & Standard convention \\
SHM $(v_0,v_E,\vesc)$ & $(238,263,544)$~km/s & Codebase production default \\
Mediator cases & Heavy ($n_\mathrm{FDM}=0$), light ($n_\mathrm{FDM}=2$) & Both scanned in the DAMIC-M PRL \\
\bottomrule
\end{tabular}
\end{center}

This reproduces the DAMIC-M PRL 2025 result \citep{DAMICM2025}: freeze-in
excluded for an ultralight mediator over 3.5--490~MeV/$c^2$, freeze-out
excluded for a heavy mediator over 2.9--21.5~MeV/$c^2$ (see
\texttt{docs/LBC\_QEDark\_Reproduction\_Guide.md} for the operational
reproduction steps).

\subsection{\SrCdSb{} (HypMat) case}
\label{sec:dme-hypmat}

No first-principles DFT wavefunctions for \SrCdSb{} are yet available
(AiiDA \texttt{bands\_workchain\_pk=9491}, pending). The DM-electron rate
for this material is therefore computed with the \textbf{scissors
approximation}: a scissor-shifted Si dielectric function stands in as a
phenomenological proxy for the true \SrCdSb{} $\varepsilon(q,\omega)$.

\paragraph{Procedure (scissors operator), step by step.}
\begin{enumerate}[leftmargin=*]
  \item Run DFT on bulk Si (PBE functional, $4\times4\times4$ $k$-grid — the
        ``fast'' template \texttt{Si\_fast\_scissor.in}), obtaining
        Kohn--Sham orbital energies and coefficients.
  \item Measure the DFT gap $\Delta_\mathrm{DFT} = \varepsilon_\mathrm{LUMO}
        - \varepsilon_\mathrm{HOMO}$ (typically 0.6--1.0~eV for PBE/Si, an
        underestimate of the true 1.12~eV Si gap — the well-known DFT
        band-gap problem).
  \item Rigidly shift all \textbf{conduction-band} energies:
  \begin{equation}
  \varepsilon_n^{(\mathrm{scissor})} = \varepsilon_n^{(\mathrm{DFT})} +
  (\text{target gap} - \Delta_\mathrm{DFT}) \qquad \forall\, n \in \text{conduction}\,,
  \end{equation}
  with target gap $=0.34$~eV (\SrCdSb{} indirect gap, working value).
  \item Recompute $\Im\,\varepsilon(q,\omega)$ in the RPA using the shifted
        energies but \textbf{unchanged} wavefunction coefficients — this
        moves \emph{where} the spectral weight sits (transition energies)
        without changing \emph{how strong} each transition is (matrix
        elements, which remain Si's).
  \item Package the result as \texttt{data/qcdark2\_epsilon/Si/Si\_fast\_gap0p34.h5},
        used identically to a native \SrCdSb{} $\varepsilon(q,\omega)$ file
        by the rate generator.
\end{enumerate}

\begin{center}
\begin{tabular}{@{}p{4.3cm}p{4cm}p{5.5cm}@{}}
\toprule
Parameter & Value & Note \\
\midrule
Band gap (DM-e threshold) & 0.34~eV (indirect, working value; R2SCAN no-SOC:
  indirect 0.560~eV / direct 0.695~eV) & Scissors target \\
Mean $e$-$h$ pair energy $\ehpair$ & \textbf{1.7778~eV} (Klein's formula,
  Section~\ref{sec:klein}) & Supersedes earlier 1.6~eV placeholder \\
Target density $\rhoT$ & \textbf{5.76~g/cm$^3$} (from $a=4.44$~\r{A},
  $c=28.196$~\r{A}, $M_W=555.96$~g/mol, $Z=3$) & Supersedes earlier
  8.0~g/cm$^3$ placeholder \\
Valence electron count & $Z_\mathrm{val}=16$ per formula unit
  (Sr~$2s$ + 2$\times$Cd~$2s$ + 2$\times$Sb~5~valence) & Used for Drude
  $\omega_p$ (Section~\ref{sec:darkphoton-hypmat}); DM-e rate keeps Si's
  4/cell (open assumption) \\
Crystal form factor source & Scissor-shifted \texttt{Si\_fast\_gap0p34.h5}
  (Si wavefunctions, shifted energies) & \textbf{Band topology and matrix
  elements are Si's}, not \SrCdSb{}'s \\
$k$-grid & $4\times4\times4$ (``fast'') & Coarser than $8\times8\times8$
  production template; noisier $\varepsilon(q,\omega)$ \\
Local field effects & Not included at this level & Most relevant at small
  $q$, i.e.\ lowest-mass DM \\
\bottomrule
\end{tabular}
\end{center}

\paragraph{Explicit residual assumptions} (full catalog in
\texttt{HypMat\_Signal\_Models\_Unified.md}, labels B1--B9): the dominant
one is that the spectral \emph{shape} of $dR/dE$ (where the peaks in
$\varepsilon_2$ sit, how fast the phase space opens above threshold)
reflects Si's band structure, not \SrCdSb{}'s. The scissors correction is
trustworthy for the \textbf{threshold location}; it is not trustworthy for
the detailed shape. This is a genuine sensitivity \textbf{projection using
a proxy form factor}, not a first-principles prediction, and must be
reported as such.
